% Clase
\documentclass[11pt,a4paper,spanish,twoside]{report}

% Órdenes auxiliares
% Español
\usepackage[spanish]{babel}
\usepackage[utf8]{inputenc}
\usepackage[T1]{fontenc}
\usepackage{lmodern}

% Imágenes
\usepackage[pdftex]{graphicx}
\usepackage{latexsym}
\usepackage{fancybox}

% Ruta para las imágenes
\graphicspath{{img/}}

% Rotaciones
\usepackage[twoside]{rotating}

% Multirow
\usepackage{multirow}

% Referencias
\usepackage[spanish]{varioref}
\usepackage[pdftex,colorlinks=true,linkcolor=black]{hyperref}

% Colores
\usepackage{color}
\usepackage{colortbl}

% Párrafos
\setlength{\parskip}{6pt}

% Code for creating empty pages
% No headers on empty pages before new chapter
\makeatletter
\def\cleardoublepage{\clearpage\if@twoside \ifodd\c@page\else
    \hbox{}
    \thispagestyle{empty}
    \newpage
    \if@twocolumn\hbox{}\newpage\fi\fi\fi}
\makeatother \clearpage{\pagestyle{empty}\cleardoublepage}

%% Entorno Listings para código fuente
\usepackage{listingsutf8}[2007/11/11]

\lstset{
  frame=Ltb, %forma del cuadro
  framerule=0pt, %ancho del borde
  aboveskip=0.5cm, %separación de los números de línea
  framexleftmargin=0.4cm, %margen externo izquierdo
  framesep=0pt,
  rulesep=.4pt,
  rulesepcolor=\color{black},
  % 
  stringstyle=\ttfamily,
  showstringspaces = false,
  basicstyle=\scriptsize,
  keywordstyle=\bfseries,
  % 
  numbers=left,
  numbersep=15pt,
  numberstyle=\tiny,
  numberfirstline= false,
  %
  inputencoding=utf8/latin1
}

% minimizar fragmentado de listados
\lstnewenvironment{listing}[1][]{
  \lstset{#1}\pagebreak[0]}{\pagebreak[0]
}


% Fancyhdr - Encabezados y pies de página
\usepackage{fancyhdr}
% Márgenes
\headsep=8mm
\footskip=14mm

% Fancy Header Style Options
\pagestyle{fancy}               % Sets fancy header and footer
\fancyfoot{}                    % Delete current footer settings

% Sin mayúsculas en la cabecera
\lhead{\nouppercase{\rightmark}}
\rhead{\nouppercase{\leftmark}}

% Capítulo
\renewcommand{\chaptermark}[1]{ % Lower Case Chapter marker style
  \markboth{\chaptername\ \thechapter.\ #1}{}} 

% Sección
\renewcommand{\sectionmark}[1]{ % Lower case Section marker style
  \markright{\thesection.\ #1}} 

% Página
\fancyhead[LE,RO]{\bfseries\thepage} % Page number (boldface) in left on even
                                     % pages and right on odd pages

% ------------------ Macro para encabezados y pies de página-------------------
%    Uso: \encabezado{pares(pag izquierda)}
% -----------------------------------------------------------------------------
\def\encabezado{
  \fancyhead[RE]{\bfseries\leftmark}     % In the right on even pages
  \fancyhead[LO]{\bfseries\rightmark}  % In the left on odd pages
  \renewcommand{\headrulewidth}{0.5pt} % Width of head rule
}
% -----------------------------------------------------------------------------

% ------------------ Macro para insertar una imagen ---------------------------
%    Uso: \imagen{nombreFichero}{Ancho(cm))}{Etiqueta}{Identificador}
% -----------------------------------------------------------------------------
\usepackage{float}
\usepackage{ifthen}
\def\imagen#1#2#3#4{
  \begin{figure}[H]
    \begin{center}
      \ifthenelse{\equal{#2}{}}
      {\includegraphics{#1}}{\resizebox{#2cm}{!}{\includegraphics{#1}}}
      \ifthenelse{\equal{#3}{}}{}{\caption{#3}}
      \label{#4}
    \end{center}
  \end{figure}
}
% -----------------------------------------------------------------------------


% ------------------ Macro para la portada ------------------------------------
%    Uso: \portada{asignatura}{titulo}{subtítulo}{autor}{fecha}
% -----------------------------------------------------------------------------
\def\portada#1#2#3#4#5{
  \thispagestyle{empty}
  \vspace*{-3.3cm}

  \begin{minipage}[t]{14cm}
    \begin{center}
      \includegraphics[scale=0.25]{logoesi}
  
      \vspace*{1.5cm}
      {\Large \textbf{Universidad de Castilla-La Mancha\\ 
          Escuela Superior de Informática}\\}
    
      \vspace*{1.2cm}
      {\huge \textbf{#1}\\}
    
      \vspace*{1.5cm}
      {\huge #2}\\{\Large #3\\}
    
      \vspace*{1.5cm}
      {\large #4\\}
      \vspace*{1.4cm}
      \large{#5}
    \end{center}
  \end{minipage}

  \newpage
%   \vspace*{1cm}
%   \thispagestyle{empty} 
%   \newpage
}
% -----------------------------------------------------------------------------


% ------------------ Macro para la licencia -----------------------------------
%    Uso: \portada{asignatura}{titulo}{subtítulo}{autor}{fecha}
% -----------------------------------------------------------------------------
\def\licencia#1{
  \thispagestyle{empty}  % Suprime la numeración de esta página
  \vspace*{16.5cm}
  \begin{small}
    \copyright~ #1. Se permite la copia, distribución y/o 
    modificación de este documento bajo los términos de la licencia de
    documentación libre GNU, versión 1.1 o cualquier versión posterior publicada
    por la {\em Free Software Foundation}, sin secciones invariantes. Puede
    consultar esta licencia en http://www.gnu.org. \\[0.2cm]
    Este documento fue compuesto con \LaTeX{}. 
  \end{small}
  \newpage
%   \thispagestyle{empty}
%   \vspace*{0cm}
%   \newpage
}
% -----------------------------------------------------------------------------




% Encabezado y pie de página
\encabezado

\setcounter{secnumdepth}{3}
\setcounter{tocdepth}{3}

\begin{document}

% Silabación extra
\hyphenation{
a-sig-na-tu-ras
au-to-ma-ti-za-rá
ca-tá-lo-go
ca-rre-ra
diag-nos-tico
in-fe-rior
man-te-ni-mien-to
per-so-nal
pro-por-cio-na-rá
pu-bli-ca-da
re-qui-si-tos
res-pecto
u-su-a-rios
vi-lla-rre-al
}


% Portada
\portada{Automatización Industrial}
{Trabajo:}{Automatización del proceso de embotellado en\\ J. García Carrión
(Daimiel)}
{Sergio de la Rubia García-Carpintero\\Miguel Millán Sánchez-Grande}
{30 de Abril de 2010}

% Licencia
\licencia{Sergio de la Rubia García-Carpintero, Miguel Millán Sánchez-Grande}

% Índices
\tableofcontents
% \listoffigures
% \listoftables

%% INICIO DEL DOCUMENTO %%%%%%%%%%%%%%%%%%%%%%%%%%%%%%%%%%%%%%%%%%%%%%%%%

\chapter{Introducción}
%soltar el rollo de xq la hemos elegido, en que consiste.
A la hora de elegir una empresa sobre la que realizar el trabajo, se buscaron
opciones cercanas geográficamente y de fácil acceso en nuestro entorno. Por
este motivo la opción elegida fue la de la empresa \textbf{J.Garcia Ca\-rrión}
(a partir de ahora \textbf{JGC}), ubicada en \textbf{Daimiel}; por tratarse de
una instalación próxima a Ciudad Real y porque el hermano de uno de los 
componentes de nuestro grupo trabaja en dicha planta.

JGC es la empresa fabricante de todos los productos \textbf{Don Simón}.
\imagen{don_simon.jpg}{4.5}{Don Simón}{D_S}
Para desarrollar su plan estratégico cuenta con distintas bodegas y 
plantas de proceso y envasado de última tecnología en distintas partes de la
geografía española, siendo una empresa pionera tanto a nivel nacional como
internacional:

\begin{itemize}
\item \textbf{Jumilla (Murcia)}, origen de la empresa, con una capacidad de 
400 millones de litros. 
\item \textbf{Daimiel (Ciudad Real)}, pionera en tecnología a nivel mundial, 
cuya alta producción de 800 millones de envases por año, la hace la más 
eficiente de Europa. Sobre ella se centra este trabajo.
\item \textbf{Gador (Almería)}, inaugurada en el año 2003, es una Planta 
especia\-li\-za\-da en tratamiento de vegetales (gazpacho, cremas de verduras
na\-tu\-ra\-les y caldos), con una capacidad de 120 millones de Kilos/año. La
tecnología de esta planta se ha proyectado con Investigación y De\-sa\-rro\-llo
propio, único en la elaboración de productos naturales vegetales. 
\item \textbf{Huelva}, planta única en el mundo por su agricultura integrada.
\end{itemize}

\section{Historia breve}
La familia JGC siempre ha estado ligada a los viñedos y a la tradición
agrícola, tiene sus orígenes en el pueblo murciano de \textbf{Jumilla}. En 
1890 la familia construyó una nueva bodega, con ciertas dimensiones para 
aquella época, debido al gran auge de la exportación del vino de Jumilla a 
Francia, esta exportación fue originada por un parásito llamado filoxera que
arrasó con los viñedos del país galo. Por tanto, se toma el año \textbf{1890} 
como él de la fundación de \textbf{GARCIA-CARRIÓN}.

La marca GARCIA-CARRIÓN, prosigue su andadura, creciendo poco a poco y
sobreponiéndose a distintas adversidades como la guerra civil. Aún así,
comenzó a distribuirse por toda España. Al continuar aumentando la demanda, 
se decide construir una nueva bodega en las proximidades de Jumilla, e
instalar el primer tren de embotellado de alta capacidad.
 
Si por algo se ha caracterizado siempre la compañía, ha sido por la innovación
y por arriesgarse a realizar cosas que otros pensaban imposibles; de
esta forma, a principios de los 80, con la implantación de las grandes
superficies en España, el envase más utilizado para el vino de mesa era la
botella de 1l retornable, lo que exigía la posesión varias plantas de
envasado distribuidas por toda España para atender la demanda nacional. Con
el fin de buscar un envase no retornable, práctico, económico y de poco peso
la compañía optó por la tecnología brik.

El lanzamiento del nuevo envase coincidió con el primer anuncio en televisión
de la compañía, que será siempre recordado por la frase : ``Voy a comer con
Don Simón''.

\chapter{Objetivo y motivación del sistema}
La apertura de la planta de Daimiel coincide con la linea general de
desarrollo e innovación que sigue la compañía JGC desde su creación, como
se ha podido observar a través de su historia. La planta objeto de estudio fue
originariamente inaugurada hace tres años, esta planta cubría una necesidad
creada a partir del aumento de la demanda de sus productos. Fue necesario 
conseguir una receta para elaborar más producto, de una forma más veloz, pero,
por supuesto, a un precio razonable. La alternativa de aumentar el imperio JGC
sobre la geografía española se planteaba como una alternativa cara de ejecutar,
así que se optó por la construcción de una sola planta en la que invertir más 
dinero para conseguir un alto rendimiento.

\imagen{planta.jpg}{9}{Planta en la actualidad}{Planta}

Pero durante la primera producción de dicha planta, un incendio la dejó hecha
cenizas. Esta planta funcionaba con máquinas de la marca \emph{Sindel}, y se 
encargaba de la producción de tinto de verano. Esa misma producción tuvo que
ser finalizada en la planta antigua que la empresa posee en Jumilla.

Por suerte, con el dinero que se consiguió con la liquidación del seguro unido
a una subvención de la Junta de Comunidades de Castilla la Mancha, se decidió
construir la planta que conocemos a día de hoy, que podemos observar en la 
figura \ref{plan} de la página \pageref{plan}. Una de las más punteras
tecnológicamente hablando, no sólo a nivel nacional, sino que también a nivel
internacional. Como posteriormente se explicará, casi todas las máquinas que
desarrollan la cadena de trabajo son de la marca alemana \textbf{KRONES}.

\imagen{plano.jpg}{5}{Plano de la planta}{plan}

\chapter{Descripción del proceso}
El proceso que se lleva a cabo en la fábrica de JGC situada en Daimiel y
sobre el cual se centra el esfuerzo de este trabajo consiste en embotellar
zumo, sangría y tinto de verano, dependendiendo de las necesidades y de la
temporada. Es decir, en general cualquier producto de Don Simón que esté 
envasado en botella de plástico.

En la \emph{figura \ref{proceso}} se representa un mapa de todo el proceso
automatizado. Como se aprecia, consta de seis etapas que se han de cumplir
secuencialmente, para que dada una preforma y el líquido que se quiere
embotellar, consigamos un palé lleno de botellas.
\imagen{proceso.pdf}{12.5}{Proceso de embotellamiento}{proceso}

\begin{itemize}
\item La primera etapa consiste en a partir de las preformas obtener botellas 
de plástico. Una preforma es, como se puede apreciar en la \emph{imagen  
\ref{preforma}}, una pequeña probeta de plástico. Esta preforma es
introducida en una máquina llamada \textbf{Sopladora}. A la salida de esta 
etapa las preformas son convertidas en botellas de plástico vacías.
\imagen{preforma-20.pdf}{4.5}{Preforma}{preforma}

\item En la siguiente etapa entra en juego no sólo la salida de la etapa
  anterior, si no que también el tinto. Ambos se compaginan en la
  \textbf{Llenadora/Taponadora}; a la salida de esta etapa se tienen las
  botellas llenas y tapadas. Aquí no se abarca, pero en esta etapa existen dos
\textbf{inspectores}: uno del tapón, que se encarga de controlar que los 
tapones están correctamente colocados; y otro de volumen, que mide el nivel de
líquido que contiene cada botella.

\item Posteriormente, se pasa a etiquetar las botellas. Para ello se 
introducen en la \textbf{Etiquetadora}. A la salida se obtienen las botellas
correctamente etiquetadas. Esta etapa también consta de un inspector para 
confirmar que las etiquetas fueron colocadas de forma correcta. Posteriormente
se profundizará más sobre ello.

\item Seguidamente, las botellas son conducidas hacia la \textbf{Agrupadora}
  que, como su nombre indica, que encarga de agrupar packs de cuatro 
bo\-te\-llas, forrándolas en plástico, de forma que su futura manipulación sea
por conjuntos, no de forma individual.

\item Después, los packs de botellas pasan a otra máquina llamada
  \textbf{Paletizadora}, que se encarga de formar capas de packs para,
posteriormente colocarlos encima de un palé.

\item Finalmente, el palé de botellas es transportado hacia la 
\textbf{Enfardadora}, donde se forra con unas capas de plástico transparente 
para evitar pérdidas en el transporte o que el producto sufra algún daño.

\end{itemize}

Cada una de las islas de automatización está controlada por un \textbf{PLC}.
Los PLC empleados son \emph{SIEMENS S7300}. Cada uno de ellos carga y guarda 
unos parámetros o ``recetas'', que son los que rigen el funcionamiento de este 
controlador, en un PC. Los PCs empleados tienen el S.O.
\emph{Windows XP (Embedded)} y un software \textbf{SCADA}, proporcionado por
el fabricante de la máquina, llamado \emph{ZENON}, desarrollado por 
\emph{COPA-DATA}. La comunicación entre PC y PLC se desarrolla por Ethernet,
aunque también puede hacerse mediante R232.

Hay que destacar que las máquinas trabajan de forma autónoma, es decir, no se
comunican unas con otras. Simplemente, mediante sensores, controlan las
entradas y las gestionan. No obstante, también hay que decir que estas máquinas
tienen la posibilidad de ser gestionadas de una forma centralizada, cosa que
es posible que se haga a corto o medio plazo en la planta.

Aparte de las diversas islas de automatización de las que se han hablado, el
transporte del producto de una de ellas a la siguiente, también está 
au\-to\-ma\-ti\-za\-do. Esto se lleva a cabo mediante diversas \textbf{cintas 
transportadoras}, mesas de acumulación y alineadoras. Pero de todo esto se 
hablará más a fondo en el capítulo siguiente. 

\chapter{Explicación de las islas de automatización}
\section{Sopladora}
La primera de nuestras islas de automatización es la \textbf{sopladora}.

Esta planta dispone de una sopladora \emph{KRONES Contiform}, con una 
capacidad de trabajo de más de 14.000 botellas a la hora.
\imagen{sopla.jpg}{8}{Sopladora}{Sopla}

Las preformas se encuentran dispuestas en una fila a la entrada de un horno. 
Cada preforma entra al túnel y se calienta por radiación. Según el tipo de 
resina de con la que está hecha y la geometría que se quiere conseguir, se 
aplica una temperatura específica (unos 110ºC de media). Para controlar la 
temperatura que se aplica a cada preforma, el horno cuenta con un 
\textbf{pirómetro}. Con la ayuda de un controlador \textbf{PID} (Proporcional
Integral Derivativo), la temperatura que se aplica a la preforma se mantiene
en todo momento constante, ya que de no ser así, podría provocarle
deformaciones.
\imagen{molde2.jpg}{5}{Molde abierto}{mol}
Una vez caliente, la preforma es conducida a un molde como el de la
figura \ref{mol}. El molde está compuesto por tres piezas que se separan
para permitir la entrada de la preforma. Una vez dentro, las piezas se unen y
desciende la sopladora sobre la boquilla de la preforma. La sopladora realiza
dos operaciones sobre la preforma: 
\begin{itemize}
\item Primero, se realiza el \textbf{estirado y presoplado}. Este consiste en
que una varilla se introduce en la preforma y la va estirando mientras la 
sopla hasta que le da la longitud final que tendrá la botella.
\item Después, se realiza el \textbf{soplado moldeante}. Este hace que la 
preforma se hinche hasta tocar las paredes del molde. Por el interior del
molde circula una corriente de agua (a una temperatura aproximada de 9ºC) que
mantiene frías las paredes. Cuando el material caliente de la preforma toca la
pared fría del molde, adquiere la forma de este y se enfría.
\end{itemize}
Para controlar los procesos de estirado y soplado, el sistema cuenta también 
con un controlador PID, que maneja el traductor de presión y la válvula 
modulante de presión.

\imagen{molde.jpg}{9}{Moldes de la sopladora}{molde}
En total, la sopladora consta de 8 moldes como se ve en la figura \ref{molde}.
Estos se encuentran dispuestos sobre una superficie giratoria que, mientras 
gira, va cogiendo las preformas que se le van sirviendo para meterlas en cada
uno de los moldes. Una vez formada cada botella, aprovechando el giro del
sistema, se vuelve a abrir el molde y va saliendo del sistema por el lado
opuesto al que entró, hacia la siguiente isla de automatización, la
\textbf{llenadora}.

\section{Llenadora/Taponadora}
Las botellas son transportadas por un carrusel hasta la llenadora.

La \textbf{llenadora/taponadora} de esta planta, es una \emph{KRONES
  Vo\-lu\-me\-tic}, como la que se muestra en la figura \ref{LLE}, con 
una capacidad de trabajo de 12.000 a 14.000 botellas a la hora, dependiendo de
la capacidad de las botellas que deben llenar.
\imagen{llenadora.jpg}{10}{Llenadora KRONES}{LLE}
Al llegar a este sistema, cada botella es enchufada por un caño. Esta máquina 
posee un total de 80 caños conectados a un depósito donde está alojado el 
líquido con el que vamos a llenar las botellas, en este caso tinto de verano.

Primero, cada caño hace una \emph{presurización} de la botella, con el 
objetivo de equilibrar la presión existente entre la botella y el caño para 
que pueda caer el líquido en la botella y que no forme espuma.

Para conocer el volumen que contiene la botella en un instante dado, cada caño 
posee un \textbf{caudalímetro}. Como el llenado de la botella se produce
mientras se transporta a la siguiente máquina del sistema, se debe configurar
correctamente la velocidad del carrusel para que le de tiempo a llenar la 
botella.

Este sistema, al igual que casi todos los sistemas de la planta, utiliza 
varias \textbf{fotocélulas de ultrasonidos} para detectar el paso de las 
botellas, y \textbf{codificadores rotatorios} (o \emph{\emph{encoder}s}) para 
calcular, mediante el conteo de pulsos, cuando deben actuar los actuadores. 
En este caso, detecta el paso de una botella y calcula el momento en que esta 
llegará al caño, para dar la orden de abrir a la válvula de llenado de dicho 
caño.

Cuando la botella sale del carrusel de llenado entra en \textbf{taponadora}
como la de la figura \ref{tap}.
\imagen{taponadora.jpg}{9}{Taponadora}{tap}
La taponadora detecta la entrada de una botella, abre una trampilla, deja caer 
el tapón y mediante un embrague, que realiza un par determinado, lo aprieta, 
tapando la botella.

A fin de controlar las botellas defectuosas, este sistema cuenta con dos 
\textbf{inspectores}:
\begin{description}
\item \textbf{Inspector de nivel}
Este primer inspector se encarga de controlar que la botella ha sido llenada 
hasta un nivel esperado. Las botellas que tienen un nivel defectuoso serán 
retiradas de la cadena. Este inspector detecta el nivel de la botella mediante 
rayos X y proporciona una salida entre 0 y 10V.

\item \textbf{Inspector de tapón}
Este inspector se encarga de controlar que la botella esté tapada 
correctamente. Para ello, se cuenta con dos \textbf{fotocélulas}. Estas 
fotocélulas se encuentran a la altura a la que debería ir la anilla del tapón.
La anilla del tapón hace más grueso el cuello de la botella. Las fotocélulas 
se sitúan a una distancia algo menor al diámetro del cuello de la botella, de
tal forma que si ambas fotocélulas detectan a la vez a la botella es que lleva
bien colocado el tapón. En cualquier otro caso, se controlará la botella para
ser retirada de la cadena.
\end{description}

Como se dijo anteriormente, este sistema controla mediante un \emph{encoder} 
las botellas defectuosas para retirarlas de la cadena. Las botellas son 
``literalmente'' empujadas fuera de la cadena, en su fase de transporte, 
mediante un \textbf{expulsador} (o \emph{pusher}). Este sistema empuja a la
botella fuera de la cinta transportadora. El \emph{encoder} cuenta los pulsos
que deben pasar desde que un inspector encuentra una botella defectuosa, hasta
que realmente esta botella pasa delante de uno de estos \emph{pusher}; es 
entonces cuando se da la orden de que este actúe.
\section{Etiquetadora}
\imagen{etiqueta.jpg}{11}{Etiquetadora}{}
Antes de llegar a la siguiente isla de automatización, encontramos una 
\textbf{mesa de acumulación}. Esto se debe a que, a partir de este nivel, cada
una de las máquinas que encontraremos es un 10\% más rápida que la anterior.
Por este motivo, para que la \textbf{etiquetadora}, nuestro siguiente 
subsistema a analizar, pueda trabajar con continuidad, es decir, sin que le 
falten los recursos que necesita para trabajar, necesita contar con un número 
considerable de botellas en stock. Cuando se alcanza este nivel, se activan 
las \textbf{cintas transportadoras} que conducen a las botellas a la 
etiquetadora. Para controlar el número de botellas almacenadas, la mesa de 
acumulación cuenta con un sensor que mide la presión que ejercen las botellas 
contra sus paredes. Cuando se alcanza la presión programada, la cinta comienza 
a moverse. Por otro lado, cuando esta presión es nula, el movimiento de la 
cinta cesa y de forma automática, también se para la etiquetadora.

Llegados a este punto, encontramos a un gran número de botellas agrupadas 
desplazándose hacia la etiquetadora, aunque, esta sólo pueda recibir las 
botellas enfiladas. Es por ello que necesitamos un mecanismo intermedio que 
las alinee. Para ello se cuenta con una \textbf{alineadora} como la de la
figura \ref{ali}. Este sistema cuenta con una serie de estrechas cintas
transportadoras, colocadas unas al lado de las otras, de forma paralela. Y
cada cinta se des\-pla\-za a una velocidad superior que la que se encuentra
inmediatamente a su izquierda. Pues bien, el sistema funciona de la siguiente
manera: las botellas llegan apiñadas por la cinta más ancha y más lenta hasta
que se encuentra en su camino con un estrechamiento. En ese momento, las
bo\-te\-llas que van primeras, son empujadas a una cinta más estrecha y más
rápida. Las bo\-te\-llas que han sido empujadas a esta cinta, se encontrarán de
nuevo con otro estrechamiento, cayendo de nuevo a otra cinta más rápida que
la anterior. El resultado será que, cuando se  llega a la última cinta, todas
las botellas que se encuentran en cabeza, estarán alineadas. 

\imagen{alineadora.jpg}{5}{Alineadora de botellas}{ali}

Por fin, las botellas llegán a la \textbf{etiquetadora}. Esta etiquetadora es 
una \emph{KRONES Contiroll}, y tiene una capacidad de trabajo de más de 14.000 
botellas a la hora.

Un tornillo sin fin va separando las botellas entrantes para evitar que lo 
hagan todas seguidas. La botella se encuentra con un primer cilindro 
giratorio, llamado \emph{de estrella}. Este recoge la botella del tornillo sin
fin y se la sirve, al otro lado, a otro cilindro giratorio en estella mucho 
mayor. Este cilindro recoge la botella en una de sus muecas y la hace también 
girar sobre sí misma. Durante su camino por el cilindro, la botella se 
encuentra en primer lugar con el sistema de etiquetado. Este, está formado por
una serie de rodillos: 
\begin{itemize}
\item Dos rodillos que contienen dos bobinas de etiquetas, una que va
sumi\-nis\-trando las etiquetas y otra que está para cuando se termine el 
primer rollo. La conmutación para cambiar entre uno y otro rollo se realiza 
automáticamente, de forma que no hace falta parar la máquina para cambiar de 
rollo.
\item El siguiente rodillo, es el de la cortadora. Antes de llevar a dicho
rodillo, tenemos una fotocélula de contraste, que detecta el borde donde se 
debe cortar la etiqueta. Esta fotocélula, coordinada con un \emph{encoder}, 
avisa de cuándo debe actuar la cortadora del rodillo sobre la etiqueta.
\item La etiqueta llega al rodillo de encolado. Este rodillo, que gira en 
sentido opuesto al anterior, recoge la etiqueta y la deja montada sobre dos
escalones, que coinciden con el inicio y el final de la etiqueta. Esto se 
hace para que cuando el cilindro con la etiqueta pase por el rodillo auxiliar
que contiene la cola, impregne la superficie de la etiqueta que se encuentre
sobre dichos escalones.
\item El último tramo del viaje de la etiqueta coincide con su encuentro con
la botella, que viene del cilindro principal, girando en sentido opuesto a
este. La cola hace que la etiqueta se pegue a la botella y el movimiento
giratorio, que la rodee por completo.
\end{itemize}

Siguiendo con el recorrido de la botella, antes de dejar el cilindro 
principal, pasa por un \textbf{inspector de etiquetas}, consistente en un
fotógrafo que compara fotos por contraste. Con ello averigua si la etiqueta
se ha colocado correctamente sobre la botella. En caso contrario, al igual
que ocurría con la llenadora/taponadora contamos con otro \emph{pusher}
que, ayudado de un \emph{encoder}, expulsará las botellas que no cumplan los
requisitos de contraste especificados.

Finalmente, la botella se encuentra con otro cilindro en estrella, que la
recoge y la deja en una cinta transportadora, donde se encuentra el 
\emph{pusher} mencionado anteriormente.

\section{Agrupadora}
A la salida de la etiquetadora encontramos otra mesa de acumulación, puesto 
que, como ocurría en el caso anterior, la siguiente isla de automatización, 
tiene la capacidad de trabajar un 10\% más rápido que la etiquetadora.

Esta isla de automatización es la \textbf{agrupadora}. La agrupadora de esta 
planta es una \emph{KRONES Variopac}.

Antes de llegar a la máquina, las botellas son ordenadas de dos en de dos y 
separadas mecánicamente en packs de cuatro como se muestra en la figura
\ref{agru}. 
\imagen{agrupa.jpg}{12}{Agrupadora de botellas}{agru}


El objetivo de la agrupadora es plastificar lotes de cuatro botellas para
luego meterlos en palés.

La agrupadora consiste, básicamente, en una cinta transportadora bajo la que 
se encuentran varios rodillos móviles y varios fijos. Los rodillos fijos se 
encargan de tensar el rollo de plástico con el que se plastificarán los lotes.
Los rodillos móviles portan el extremo del rollo y se mueven describiendo 
arcos sobre los packs, dejándolos cubiertos.

Una vez cubiertos, los packs pasan por un \textbf{túnel de retractilado}, 
donde se calientan a una temperatura de unos 235ºC, haciendo que el plástico
envolvente se encoja y fije las botellas del pack.

\section{Paletizadora}
El último paso antes de obtener un palé, es formar las capas de packs que lo
componen y apilarlas.

En este caso, la paletizadora no es \emph{KRONES}, como venía ocurriendo con 
las islas de automatización anteriores, sino \emph{ELETTRIC80}, una firma
italiana.

Como puede observarse en la figura \ref{pal}, la paletizadora consiste de: una
cinta transportadora, por la que circulan los packs; y de tres robots de tipo
\textbf{angular}, que se encargan de manipular los packs para formar capas,
poder dichas capas en un palé y separar cada capa del palé con un cartón,
respectivamente.

\imagen{pale.jpg}{12}{Paletizadora}{pal}

Vamos a ver con detalle como trabajan cada uno de estos tres robots:
\begin{itemize}
\item El primer robot es un \emph{FANUC M410i}, tiene una capacidad de carga de
50kg y es el encargado de formar las capas del palé. Tiene un actuador,
consistente en una pinza, que se encarga de coger y soltar los packs en una
determinada posición y orientación en la misma cinta transportadora. De esta
forma, cuando los packs llegan al final de la cinta, llegan colocados para 
formar un mosaico con la forma de un palé, como si se tratase de un tetris de
packs. Gracias a unos sensores colocados en la cinta y a una memoria, el
controlador del robot conoce en cada momento donde están los packs que tiene
que colocar y donde y como tiene que colocarlos.
\item El segundo robot es un \emph{FANUC M710i}, tiene una capacidad de carga
de 300kg y es el encargado de coger las capas formadas por el robot anterior y
apilarlas sobre un palé. Este robot se limita a desplazar su actuador cuando
un final de carrera, colocado al final de la cinta, detecta la llegada de una
capa de packs. El actuador consiste en una pinza que coge la capa y la
desplaza hasta una plataforma cercana donde se encuentra el palé.
\item El tercer robot es del mismo modelo que el anterior. Realiza un trabajo
bastante sincronizado con este, ya que se encarga de coger un carton y
colocarlo sobre una capa, cuando el segundo robot ha colocado una capa de
packs. Normalmente los palés constan de pilas de cuatro capas, para las
botellas de 2l; o de pilas de seis capas, para las botellas de 1'5l. Cuando
se ha completado el palé, la plataforma con cadenas sobre la que se
encuentra, se  lo lleva. En ese momento, la labor de este robot será la de,
en vez de colocar un nuevo cartón, coloca un nuevo palé sobre la
plataforma. La labor del segundo robot, sin embargo, no se ve alterada, ya
que su labor continúa siendo la de coger capas de packs y colocarlos sobre
los palés directamente, o sobre los cartones que tapan una capa anteriormente
colocada. 
\end{itemize}
Tanto la cinta transportadora, como los robots, se encuentran conectados a un
concentrador de red con el PLC que los controla. Esto posibilita llevar un
control coordinado de todo el proceso.

Los palés completados son recogidos por unos \textbf{vehículos guiados 
automáticamente} (o \emph{AGVs}), similares al que podemos ver en la figura
\ref{agv}, para ser transportados hacia nuestra siguiente
y última isla de automatización, la \textbf{enfardadora}.
\imagen{agv.jpg}{10}{Vehículo guiado automáticamente}{agv}

\section{Enfardadora}
La tarea de la \textbf{enfardadora} consiste en forrar de plástico los 
laterales de los palés que le van llegando. Con ello se consigue la 
consistencia y robustez del palé, clave a la hora del proceso de transporte.

Al igual que la anterior, este sistema tampoco es \emph{KRONES}, sino 
\emph{ELETTRIC80}.

\imagen{enfar.jpg}{10}{Enfardadora de palés}{}

La enfardadora tiene tres partes bien diferenciadas:
\begin{description}
\item \textbf{Mesa de rotación}
La mesa de rotación es una sección de la propia cinta transportadora que tiene 
la funcionalidad de girar en un plano paralelo al del suelo, pero sobre un eje 
perpendicular a este. Esta característica se aprovecha para, dado un rollo de 
plástico colocado cerca del palé, ir enrollando con el plástico el contenido 
del palé varias veces, de arriba a abajo y viceversa, según las 
características del palé y de su contenido.
\item \textbf{Pisón}
Una vez colocado el palé encima de la mesa de rotación, se fija al suelo con 
ayuda del pisón. El pilón es una columna anclada al suelo que posee, en su 
parte superior, una plataforma móvil que desciende hasta aprisionar contra la
mesa de rotación al palé. Dicha plataforma gira junto con el giro de la mesa de
rotación. Además, el pisón tiene la funcionalidad de medir la altura del palé 
para decírselo al robot. Esto lo hace ayudado por una resistencia lineal
variable; el valor de la resistencia determinará la altura del palé.
\item \textbf{Robot}
La enfardadora posee además un robot de tipo \textbf{angular}. Este es el
encargado de colocar el rollo de plástico de tal forma que, el movimiento de la
mesa de rotación, vaya enrollando el plástico en el palé. El movimiento del
brazo robótico posibilita forrar el palé de arriba a abajo con un simple
movimiento ascendente o descendente del portador del rollo. Pero además, este
modelo, también permite cambiar de rollo de enfardado sin tener que parar la
máquina. El robot posee dos rollos de plástico; uno es el que se está
utilizando para realizar el enfardado y el otro es el de repuesto, cuando se
acaba el rollo, el propio robot desecha el rollo gastado y acude a su
portarrollos a por un recambio para poder continuar su labor. El operario por
su parte puede cambiar en ese momento el rollo vacío del portarrollos para
cuando vuelva a gastarse el rollo actual.
\end{description}

La enfardadora maneja una serie de variables de actuación según la
na\-tu\-ra\-le\-za del palé a enfardar. Así se podrá configurar la velocidad de
rotación de la mesa y la tensión ejercida por el robot a la hora de colocar el
rollo de plástico para el forrado del palé.

Los palés forrados continúan por la cinta transportadora contigua hasta el 
extremo de esta. Una vez allí, una de los numerosos vehículos con control 
automatizado se encargará de recoger el palé y llevarlo al almacen.

\chapter{Beneficios de la automatización}
Introducir nuevas tecnologías en una empresa tradicional como JGC siempre
resulta un proceso complicado y que no agrada a todo el mundo. Pero no se
puede negar que supuso un empujón muy importante tanto a nivel económico como
a nivel de producción para la empresa. Errónamente se asocia el concepto de
automatización con una solución para reducir la cantidad de empleados en
plantilla, bien esta empresa es un ejemplo de que esto no sucede, ya que al
mecanizar el proceso relacionado con el empaquetado y relleno de botellas
ayudo a la expansión de la empresa creando más puesto de trabajos, eso sí
distintos a los que había con anterioridad. Aparte del beneficio económico de
reducir el número de empleados de ese área hubo otros beneficios adicionales
mayores:

\begin{description}
\item \textbf{Aumento de la eficiencia}

Los costos de producción se redujeron drásticamente al aumentar las
unidades de producto fabricadas por unidad de tiempo. 

\item \textbf{Incremento del volumen de producción}

Al aumentar el número de unidades fabricadas por unidad de tiempo se aumentó
la cantidad de unidades producidas así como el número de clientes a los que
poder atender.

\item \textbf{Estandarización de los procesos}

De esta forma se logró que los productos tuvieran siempre las mismas
características, al hacer el proceso repetitivo y siguiendo los mismos
pasos. Se puede tener la certeza que dos vasos de cualquier producto Don
Simón van a tener siempre el mismo sabor, color, densidad, etc.

\item \textbf{Reducción de los problemas de calidad}

Consecuentemente con la estandarización de procesos se consiguió aumentar la
calidad, se eliminó cualquier error posible relacionado con un despiste
humano, ya fuera debido al cansancio o a una negligencia. 

\item \textbf{Aumento de la competitividad}

Todas estas mejoras, más y mejor producto en menos cantidad de
tiempo, tienen una clara repercusión en la competitividad de la empresa,
haciendo más fácil el cubrir diversas áreas de comercio. 

\item \textbf{Centralización de producción}

JGC es una de las empresas más grandes e importantes de España en el sector,
y como se ha hablado en la introducción consta solo de cuatro fabricas. Esto
es debido a la gran productividad asociada al uso de la automatización, no es
necesaria la construcción de mas sedes para cubrir las necesidad

\end{description}

Como demostración práctica de la utilidad de la automatización, la empresa
esta haciendo frente a la crisis que experimentamos en estos días con un
rebaja considerable en los precios de sus productos. Para realizar esto, JGC,
se está preparando para multiplicar la capacidad de producción 1,5 veces en
2014. Con el incremento de la producción pretende atender la demanda en
mercados internacionales para elevar su porcentaje de ventas al exterior
desde el 35\% actual al 60\%. Aún con la crisis la empresa tenía previsto
cerrar el ejercicio en curso con un beneficio antes de impuestos 22 millones
de euros, lo que supone un incremento del 22\% respecto al año anterior, y
elevar sus ventas un 12\%, hasta alcanzar los 650 millones de euros, en un
año en el que ha recortado sus precios un 20\% de media para hacer frente al
descenso provocado por la anteriormente nombrada crisis económica.

\begin{thebibliography}{99}
\bibitem{JGC} Página web de J.García Carrión: \url{http://www.donsimon.com}
\bibitem{Kro} Página web de KRONES: \url{http://www.krones.de}
\bibitem{You} Consulta de videos de máquinas:
  \url{http://www.youtube.com}
\bibitem{Ele} Página web de Elettric80:  \url{http://www.elettric80.it/}
\end{thebibliography}
%\bibliographystyle{plain} 
%\bibliography{t2}

\end{document}
